\chapter{总结与未来展望}
\section{研究总结}
本文以PW对连接成对搬运和腔室周期，以Pjob库存记录产品来源，建立完工时间优先的两阶段MILP。第一阶段联合确定配对、分腔、周期和资源次序，第二阶段保持产品完工上界及主要结构，优化等待与节拍。配置E的30组阶段内对照均保持产品完工并减少路线等待；配置D的单产品连续加工与清洁实例中，总等待由642降至54~s。权重对照中，16片实例以增加16~s等待换取8~s节拍偏差下降，24片实例保持不变，体现了两类时序指标的取舍。

结构化初始解改善了限时可行上界。四实例168次正常分离对照中，两个混合实例在三个训练种子下均从SPBS获得改善；新增混合场景S2的平均产品完工经初始化由521降至492~s。共享SPBS输入后的六个2片实例、414次算法比较中，完整选割组合相对HEM加SPBS的实例平均配对改善为0.0776\%，胜、平、负为2、4、0。完整组合在两例上优于三个组件消融，其余四例持平；根分离下贪心解码优于束搜索的结果说明解码收益随设置变化。

有限辅助库存影响连续加工的完工时间。单产品$2\times2$场景中，库存由4片增至8片使已证最优值由309降至303~s，$4\times1$保持296~s；四片辅助库存可通过重复借用完成连续加工与满片清洁。身份解码、完整模型检查与事件重放验证了324~s的2片最优排程和920~s的16片可行排程，并分别给出产品回片和全部占用结束时刻。

\section{产品完工瓶颈改进研究的结论}
% 该独立研究尚未完成，保留标题与空间，实验验收后填写。
\vspace*{0.8cm}
\section{未来展望}
后续将比较等总求解预算下的目标组织与较大实例组件效果，将构造、构模、SPBS、学习搜索及验证纳入统一截止时间。多产品长链和100片排程还需可扩展的构造与修复方法；产品完工瓶颈、解池共识与分歧、耦合邻域搜索需继续完成实现和消融。评价保留失败、无改善及资源停止记录。
